\documentclass[conference]{IEEEtran}

\usepackage[T1]{fontenc}
\usepackage[utf8]{inputenc}
\usepackage{cite}
\usepackage{amsmath}
\usepackage{booktabs}
\usepackage{multirow}
\usepackage{graphicx}
\usepackage{textcomp}
\usepackage{xcolor}

% TODO(Prabin): fill in the real title, authors, and affiliations before
% submission. These are placeholders so the document compiles.
\title{Closing the Floor-Contact Gap in General Motion Retargeting:\\
A Contact-Aware Correction Layer for Humanoid Motion Retargeting}

\author{
  \IEEEauthorblockN{Author Name(s) --- TODO}
  \IEEEauthorblockA{Affiliation --- TODO \\
  email@example.com --- TODO}
}

\begin{document}

\maketitle

\begin{abstract}
General Motion Retargeting (GMR) has become a standard tool for mapping
human motion-capture data onto humanoid robots for imitation-learning
pipelines, but its own benchmark explicitly excludes motions with complex
floor contact --- crawling, falling, prone/supine lying, and getting up.
Applied out-of-box to exactly this excluded class, GMR produces 13--16~cm
of floor penetration on up to 91\% of frames, an order of magnitude worse
than its own clean-locomotion clips (1~cm, 0.3\%). Floor contact is common
in fall recovery, manipulation from the ground, and full-body interaction
--- all targets for humanoid deployment --- so this is not a corner case.

We present a contact-aware correction layer that sits on top of GMR's own
solver without modifying it, and evaluate it with a reference-free
physics-plausibility suite built for this purpose (no ground-truth robot
motion required). On the full 77-clip LAFAN1 corpus (34 floor-contact / 43
locomotion clips), our method eliminates floor penetration and
self-collision almost entirely relative to GMR's own published height-fix
baseline (floor class: penetration 2.76~cm~$\to$~0.00~cm, self-collision
6.3\%~$\to$~0.002\%) and raises an un-gameable joint contact-correctness
metric from 0.2\% to 98.8\% of frames --- a metric we designed specifically
because a naive constant-offset baseline can satisfy simpler, single-axis
versions of the same check while leaving the robot floating well above the
ground at held frames.

We report two instantiations of this correction layer. A multi-stage,
per-frame pipeline achieves the strongest contact precision but carries a
9.8--15.5\% peak-joint-velocity cost relative to the height-fix baseline,
traced to an architectural property of independent per-frame contact
correction: adjacent, nearly-identical frames can select different,
equally-valid null-space solutions, producing joint-velocity artifacts that
aggregate contact metrics do not surface. A global, whole-trajectory
reformulation removes this artifact by construction and reduces peak joint
velocity and jerk several-fold below both baselines, at a currently lower
level of contact precision (evaluated on 36 of 77 corpus clips at time of
writing). We report both pipelines and their trade-off rather than
declaring one solved.

The kinematic result is necessary but not sufficient for downstream use;
imitation-policy training on the retargeted corpus is in progress and its
results, which will make the precision/smoothness trade-off decidable
against real downstream performance, are planned as part of this paper.
\end{abstract}

\begin{IEEEkeywords}
humanoid robots, motion retargeting, floor contact, imitation learning, contact-aware control
\end{IEEEkeywords}

\section{Introduction}
\label{sec:introduction}

\subsection{Motivation}
\label{sec:motivation}

Training humanoid whole-body control policies via imitation learning
requires large volumes of robot-executable reference motion. Since motion
capture is collected on humans, every such pipeline needs a retargeting step
that maps human kinematics onto a specific robot's morphology while
respecting its joint limits, self-collision geometry, and physical
plausibility. General Motion Retargeting (GMR) has emerged as a widely used
tool for this step: a per-frame inverse-kinematics solver (QP-based, built on
\texttt{mink}) that tracks human landmark targets on the robot subject to
hard joint limits.

GMR's own published evaluation, however, explicitly excludes motions with
complex floor contact: the authors state that they ``do not include motions
with complex interaction with the environment, such as crawling or getting
up from the floor''~\cite{araujo2025gmr}. This is a deliberate, documented
scoping choice, not an oversight. But it leaves an open question with direct
practical consequences: humanoid robots that fall, or that need to
manipulate objects from the ground, or that recover from being knocked over,
need exactly this class of motion as training data. If the standard
retargeting tool cannot handle it, either the training pipeline silently
drops this data, or it trains on badly-retargeted data without knowing it.

This is not a narrow gap. A broader survey of the human-to-humanoid
retargeting literature (contact-aware sequence retargeters, kinodynamic
refiners, learning-based retargeters --- Section~\ref{sec:related-work})
finds the same exclusion pattern repeated: contact modeling in this line of
work stops at feet and hands, target-side feasibility editing addresses foot
placement and root trajectories but not knee, pelvis, or forearm support,
and every hardware demonstration of a real floor transition to date treats
it as an \textit{emergent behavior of a learned policy} rather than the
output of a retargeter. We share this observation with a larger,
forward-looking research plan for this line of work (\texttt{paperIdea3.md},
in this repository) that scopes a full contact-template, hardware-validated
system as a 2027 target. The present paper is deliberately narrower: it
validates the first, load-bearing claim that plan depends on --- that GMR's
exclusion is a real, measurable gap, and that a targeted contact-aware layer
on top of its own solver closes most of it kinematically --- before
committing to the larger system. We return to that larger scope in
Section~\ref{sec:further-future-work}.

\subsection{Related Work}
\label{sec:related-work}

\noindent\textbf{Per-frame kinematic and teleoperation retargeters.}
\cite{darvish2019geometric} established geometric per-frame IK retargeting
validated across multiple humanoid platforms and remains a widely
referenced baseline pattern in this literature. Such systems enforce joint
limits and simple center-of-mass or zero-moment-point constraints per frame,
but have no explicit multi-link floor-contact modeling and no whole-clip
temporal optimization --- smoothness is whatever the per-frame solve and a
downstream controller happen to produce.

\noindent\textbf{Offline whole-sequence optimization retargeters.}
\cite{ayusawa2017morphing} solve human-model identification,
human-to-robot morphing, and robot motion planning as one simultaneous
trajectory optimization, enforcing joint limits and a simplified
linear-inverted-pendulum center-of-mass/zero-moment-point balance
constraint. This is the same style of quasi-static balance check an earlier
phase of this project built and deliberately left disabled by default
(Section~\ref{sec:limitations}), because a kinematic center-of-mass check
cannot distinguish genuine imbalance from momentum-carried dynamic posture
--- a limitation of this entire class of check, not particular to our own
attempt at it. \cite{otani2017adaptive}, building on
\cite{difava2016multicontact}, extend this line to explicit multi-contact
retargeting (hands, feet, and environment contacts jointly) and remain a
widely-cited foundation for later contact-aware pipelines. All three are,
like ours, target-side geometric or optimization-based methods rather than
RL- or controller-side ones; none evaluate on floor-contact postural
transitions (kneeling, prone/supine, crawling, getting up) as a first-class
case, and each assumes the robot stays continuously balanced and in known
contact with the ground --- an assumption every clip in our floor-contact
class violates by construction.

\noindent\textbf{Modern contact-aware sequence retargeters.}
\cite{jeong2025rigunification} unifies diverse motion sources onto a
canonical rig, then refines trajectories to enforce foot-contact constraints
and stability, validated on 12 simulated and 3 real humanoids --- the most
recent, and methodologically closest, prior work to ours. It is explicitly
scoped to foot-fixed, upper-body-expressive motion: its own feasibility
constraint requires both feet to remain in continuous ground contact
throughout, which structurally excludes exactly the motion class this paper
targets, where the robot must lift, reposition, and re-plant feet, or bear
weight through knees, pelvis, or a prone torso. \cite{yang2025omniretarget}
(OmniRetarget) preserves human-environment-object contact relationships
through an interaction mesh and Laplacian deformation, and is highly
influential as a data-generation backbone for later reinforcement-learning
pipelines; its contact modeling, like ours, is target-side and geometric,
but its demonstrated scope is loco-manipulation and terrain interaction, not
floor-contact postural transitions. \cite{jeong2025core} (CoRe) is
philosophically closest to our own position --- fix the geometry before
anything downstream has to compensate for it --- applying contact-aware
optimization refinement before an RL policy, but it too is scoped to
loco-manipulation contact (foot sliding, floating), not knee, pelvis, or
prone floor support.

\noindent\textbf{Retargeting as a reinforcement-learning data-generation
front end.} A large and fast-growing body of recent work --- TeleGate
\cite{li2026telegate}, the Humanoid Manipulation Interface
\cite{nai2026humi}, HumanPlus \cite{fu2024humanplus}, and the tracking-policy
ecosystem those and GMR feed, including BeyondMimic
\cite{liao2025beyondmimic} --- demonstrates that floor-contact transitions
(standing up, fall recovery, kneeling) are achievable on real hardware
today. In every one of these systems the transition is a \textit{learned
policy behavior}, executed from a reference trajectory the retargeter did
not need to get right at the contact level; the retargeter's job is only to
produce a plausible-enough reference for the policy to correct. Our
contribution is upstream of and complementary to this line of work: we ask
whether the retargeted reference itself can be made contact-correct by
construction, kinematically, before any RL policy sees it, matching this
literature's own standard practice of validating reference-motion quality
as a step distinct from policy training (Section~\ref{sec:scope}) --- and
directly motivating the policy-training evaluation planned for this paper
(Section~\ref{sec:status-next-steps}).

\noindent\textbf{A kinematic-only correction is a deliberate scope choice,
not an oversight.} An alternative to our approach is to add dynamics
directly to the correction layer: DynaRetarget
\cite{dhedin2026dynaretarget} uses sampling-based trajectory optimization to
convert imperfect kinematic retargets into dynamically feasible motions for
loco-manipulation. We do not take that route here, so that the correction
layer stays solver-agnostic and inexpensive enough to run offline over a
full 77-clip corpus without kinodynamic optimization or a physics simulator
in the loop (Section~\ref{sec:overview}). Whether the resulting kinematic,
geometric contact-correctness guarantee is also sufficient for dynamic
trackability is exactly the open question the planned policy-training
evaluation (Section~\ref{sec:status-next-steps}) answers --- this paper's
contribution stops at the kinematic guarantee, deliberately.

\noindent\textbf{Where this leaves the gap.} Across all four groups above,
floor-contact postural transitions are either out of scope by explicit
assumption (continuous ground contact, feet fixed), addressed only for
feet, hands, or manipulated objects rather than knees, pelvis, or a prone
torso, or resolved at the policy level rather than the retargeter level. We
are not aware of prior work that treats crawling, kneeling, prone/supine
lying, or getting up as a first-class kinematic retargeting target with an
explicit, non-gameable contact-correctness guarantee
(Section~\ref{sec:grounding}). That is the gap this paper closes.

\subsection{The Gap}
\label{sec:gap}

We first quantify the gap. Running GMR out-of-box on LAFAN1 clips it did not
validate against --- clips involving crawling, falling, and getting up ---
we measure 13--16~cm of floor penetration on up to 91\% of frames, an
order-of-magnitude worse than GMR's own clean-locomotion clips (1~cm, 0.3\%
of frames). This is confirmed at full-corpus scale (77 LAFAN1 clips, split
34 floor-contact / 43 locomotion by a multi-surface human contact detector,
not just a hip-height heuristic --- hip-height alone undercounts floor
contact by missing hand/knee-supported poses during brief stumbles).

We also show that GMR's own published mitigation --- a single clip-wide
constant height offset --- cannot fix this: applying it removes most of the
penetration but leaves the robot floating 6--18~cm above the floor on
average during exactly the frames that should show tight ground contact,
because a single rigid shift cannot simultaneously satisfy a clip's standing
phase and its lying phase. We built a specific metric to make this failure
mode visible and non-gameable (Section~\ref{sec:grounding}), since floor
penetration alone and floating alone can each be minimized independently by
a shift in the wrong direction.

\subsection{Contributions}
\label{sec:contributions}

\begin{enumerate}
  \item \textbf{A contact-aware correction layer}, in two instantiations,
    that sits entirely on top of GMR's own solver with zero changes to
    GMR's core tracking logic: a multi-stage per-frame pipeline (floor and
    self-collision correction, held-aware temporal smoothing, rate-limited
    correction) and a global whole-trajectory pipeline (a single sparse QP
    jointly enforcing floor/collision correctness, contact anchoring, and
    temporal smoothness), both terminating in a local grounding envelope
    with an algebraic zero-penetration guarantee.
  \item \textbf{A reference-free physics-plausibility evaluation suite} ---
    no motion-capture ground truth on the robot is available, so every
    metric we report (floor penetration, self-collision, contact
    correctness, tracking fidelity, smoothness) is computed from the
    retargeted motion and the robot's own geometry alone.
  \item \textbf{A non-gameable joint contact-correctness metric}
    (Section~\ref{sec:grounding}): a naive constant Z-shift can
    independently minimize floor penetration or eliminate floating, but not
    both at the same held frame simultaneously --- our metric requires
    both, closing a gaming loophole present in simpler, single-axis
    versions of the same check.
  \item \textbf{Full 77-clip LAFAN1 corpus validation}, floor/locomotion
    class split, against GMR's own published height-fix baseline.
  \item \textbf{An honest architectural finding and open trade-off}:
    per-frame independent contact correction can select different,
    equally-valid null-space solutions on adjacent frames, producing a
    visible ``branch-flip'' artifact that aggregate corpus metrics do not
    surface --- caught only by watching renders. We trace this to the
    correction's lack of temporal coupling, characterize a partial
    per-frame mitigation, and report results from an alternative
    whole-trajectory optimization formulation that removes the artifact
    structurally but currently trails the per-frame approach on contact
    precision. We present this as an open problem, not a solved one.
\end{enumerate}

\subsection{Scope}
\label{sec:scope}

This is a module/polish paper, not a new solver. We do not replace GMR's own
tracking IK; we validate that a targeted contact-correctness layer on top of
it closes a real, documented gap in its applicability, and we report the
result honestly including where it does not yet fully close. Kinematic
plausibility --- physically executable, contact-correct, smooth motion ---
is necessary but not sufficient for downstream use, consistent with the
standard practice in this literature of validating retargeting quality
before the separate, more expensive step of training a tracking policy on
it. Imitation-policy training on the retargeted corpus, which closes that
gap, is in progress at the time of writing and its results are planned as
part of this paper (Section~\ref{sec:status-next-steps}).

\section{Methods}
\label{sec:methods}

\subsection{Overview}
\label{sec:overview}

Our pipeline consumes a human motion-capture clip and GMR's own per-frame
retarget of it, and produces a contact-corrected, smoothed, floor-safe
trajectory for the Unitree G1 humanoid. Every stage after GMR's own solve is
ours; GMR itself (its tracking cost, its joint-limit handling, its solver)
is used completely unmodified. We present two instantiations of the
correction layer that sits on top of GMR's output: a \textbf{multi-stage}
pipeline built from a sequence of per-frame corrections
(Section~\ref{sec:multistage}), and a \textbf{global} pipeline that replaces
that sequence with a single whole-trajectory optimization
(Section~\ref{sec:global}). Both share the same canonical-human contact
labeling (Section~\ref{sec:canonical}), the same evaluation protocol
(Section~\ref{sec:grounding}), and the same final grounding step.
Section~\ref{sec:comparing} compares them.

Throughout this paper, ``contact-aware'' and ``contact-first'' refer to
\textit{geometric} contact handling: per-frame detection of which effector
is in contact (Section~\ref{sec:canonical}), and enforcement of
floor-clearance, held-position, and self-collision constraints derived from
that detection. Neither pipeline models contact forces, friction, or joint
torque --- the dynamics of contact are outside this paper's scope by design
(Sections~\ref{sec:introduction} and~\ref{sec:related-work}), the same
scope every kinematic retargeter we compare against shares.
\cite{ayusawa2017morphing}'s simplified zero-moment-point check is the
closest any of them come to a dynamics proxy, and even it assumes
quasi-static balance rather than modeling contact forces directly. The
planned policy-training evaluation (Section~\ref{sec:status-next-steps}) is
the step that tests whether a geometrically contact-correct reference is
also dynamically trackable; this paper's contribution stops at the
kinematic guarantee.

\subsection{Canonical Human and Contact Detection}
\label{sec:canonical}

Raw BVH motion capture (LAFAN1, 30~fps) is mapped to a robot-agnostic
canonical skeleton: per-frame landmark positions and orientations derived
from bone geometry, with hand orientation taken directly from GMR's own
raw-bone-rotation signal to preserve wrist twist. Contact is a per-frame
detection gate, not a correction: an effector is flagged ``in contact'' only
when it is simultaneously below a height threshold and moving slower than
0.4~m/s (7~cm for feet, 8~cm for hands) --- height alone is insufficient,
since a fast-swinging limb can pass close to the floor without being
planted. Crossing the threshold marks a frame as a candidate for anchoring
by a later stage; it does not move anything by itself. This detector covers
feet and hands and is shared verbatim by every downstream stage that needs
to know which frames are in contact.

For corpus-level evaluation, clips are additionally classified into a
floor-contact class and a locomotion class using a broader six-region human
contact labeler (feet, hands, knees, elbows, pelvis, torso; 5/8/15~cm
respectively), since a feet-only or hip-height signal undercounts floor
contact on clips where support is briefly taken through a hand or knee.

\subsection{Multi-Stage Pipeline}
\label{sec:multistage}

\noindent\textbf{Floor and self-collision clamp.} For each frame, a
two-phase sequential damped-least-squares correction runs per limb chain
(hip$\to$knee$\to$ankle, shoulder$\to$elbow$\to$wrist), proximal-to-distal.
Phase~1 resolves floor violations and pins held effectors to their contact
target; Phase~2, run strictly afterward rather than jointly with Phase~1,
resolves self-collision using the same contact-normal and
relative-Jacobian formulation as the collision term in the global pipeline
(Section~\ref{sec:global}).

\noindent\textbf{Held-aware smoothing.} A closed-form per-joint tridiagonal
smoother removes joint-velocity spikes. Held-effector degrees of freedom are
locked to their input value at high weight during contact, ramped in and
out over five frames, so smoothing does not erode a contact correction the
clamp stage already made. The floor/self-collision clamp is re-applied once
more after smoothing, since geometry-blind smoothing can reintroduce a
violation at a non-held frame.

\noindent\textbf{Rate-limited correction.} A temporal trust region on the
applied correction bounds how fast it can change from one frame to the
next, independent of the pose itself. This mitigates branch-flip artifacts
(discrete jumps between distinct, equally-valid per-frame solutions on a
redundant kinematic chain --- Section~\ref{sec:disc-tradeoff}) without
addressing their cause.

\noindent\textbf{Grounding.} See Section~\ref{sec:grounding} below (shared
with the global pipeline).

\subsection{Global Pipeline}
\label{sec:global}

The global pipeline replaces the clamp / smoothing / rate-limit sequence in
Section~\ref{sec:multistage} with one whole-trajectory optimization. GMR's
raw per-frame output is consumed directly as input.

\noindent\textbf{Floor and self-collision correction.} A sparse QP over the
actuated joint trajectory of the entire clip, solved by sequential convex
approximation: at each outer iteration, floor and self-collision
constraints are linearized into soft, slack-penalized rows (never
primal-infeasible), a trust region bounds the step, and a smoothness term
penalizes the change in the SCA correction between adjacent frames, in the
same objective as the floor, collision, and contact-tracking terms. Held
effectors are anchored to a per-region target: contact timing comes from
the canonical human's own labels (Section~\ref{sec:canonical}), a stillness
sub-check on the robot's own trajectory splits each labeled contact
interval into stationary sub-runs (anchored to a single position, ramped
in/out over five frames) and repositioning sub-frames (lightly regularized
toward their own per-frame position), and the anchor's target height is a
fixed per-effector geometric constant (the robot's own origin-to-support-point
offset at a neutral pose), not a value computed from the possibly-unreliable
input trajectory. Floor penetration is intentionally excluded from this
stage's own convergence scoring, since its solve variable is the robot's
joints only --- a pelvis-level violation on a floor-contact pose is outside
what this stage can correct, and grounding (Section~\ref{sec:grounding})
resolves it regardless of this stage's outcome; scoring against it would
only reject otherwise-improving self-collision corrections. A
cross-iteration keep-best selection guarantees the result is never worse
than GMR's own raw output on this stage's own objective.

\noindent\textbf{Smoothing and grounding.} The corrected trajectory passes
through the same held-aware smoothing stage used in
Section~\ref{sec:multistage}, then grounding (Section~\ref{sec:grounding}).
No per-frame re-clamp or rate limiter is applied --- the whole-trajectory
formulation's own smoothness term is the temporal-coupling mechanism, in
place of the multi-stage pipeline's rate limiter.

Because its smoothness term shares an objective with its floor, collision,
and tracking terms, a solution that alternates between distinct null-space
branches on adjacent frames costs the objective directly; the global
pipeline cannot produce a branch-flip by construction, independent of any
per-frame tuning.

\subsection{Grounding and Evaluation}
\label{sec:grounding}

\noindent\textbf{Grounding.} A final, deterministic pass computes the
per-frame vertical shift required to clear the floor, applies a causal
max-filter and Gaussian smoothing to that requirement, and takes the
pointwise maximum against the raw requirement. This guarantees the shifted
trajectory has zero floor penetration by construction. Unlike GMR's own
height-fix --- one rigid per-clip shift --- this is computed per frame, so
it does not force a single Z-offset to serve both a clip's standing phase
and its lying phase.

\noindent\textbf{Metrics.} No ground-truth robot motion exists, so every
metric is computed from the output trajectory and the robot's own vetted
collision geometry: floor penetration and self-collision incidence
(mesh-exact lowest point vs.\ $z=0$, k-hop-filtered for anatomically
adjacent pairs); worst float (height of a held effector's support point
above the floor); tracking fidelity (FK'd body position/orientation vs.\
GMR's own scaled human target); jerk, peak joint velocity, and
velocity-spike count; and skate (horizontal drift of a held contact point
from its own contact-onset position).

\noindent\textbf{Joint contact-correctness.} A held frame passes only if
its support point is simultaneously within a tight band of the floor and
the whole-body mesh clears the floor everywhere at that same frame. Both
conditions are necessary: a metric that checks either alone is satisfiable
by a per-clip constant Z-shift with no notion of per-frame contact at all,
since a rigid shift can be tuned to land in the right place at a clip's
labeled held frames while leaving the rest of the clip penetrating or
floating arbitrarily. Requiring both simultaneously is not satisfiable by
any single rigid shift on a multi-phase clip
(Section~\ref{sec:results-oracle}), which is why we report this joint
metric as the primary measure of contact quality rather than penetration or
float in isolation.

\subsection{Comparing the Two Pipelines}
\label{sec:comparing}

The multi-stage pipeline (Section~\ref{sec:multistage}) benefits from a
large total per-frame solve budget --- a full iterative local correction at
every frame --- at the cost of no explicit temporal coupling between
frames, which manifests as branch-flip artifacts and their associated
joint-velocity cost. The global pipeline (Section~\ref{sec:global}) removes
that artifact by construction but currently allocates a smaller total solve
budget (a bounded number of whole-clip outer iterations) to contact
precision. Section~\ref{sec:results} reports both.

\section{Results}
\label{sec:results}

All numbers are from the 77-clip LAFAN1 corpus (34 floor-contact / 43
locomotion clips) evaluated with the reference-free metric suite
(Section~\ref{sec:grounding}); the multi-stage pipeline
(Section~\ref{sec:results-multistage}) and \texttt{gmr\_heightfix}
(Section~\ref{sec:results-heightfix}) are evaluated on the complete corpus,
the global pipeline (Section~\ref{sec:results-global}) on the 36/77 clips
completed at time of writing. All comparisons are against
\texttt{gmr\_heightfix} --- GMR's own published constant-offset mitigation
--- not unmodified GMR output, which is strictly worse and reported only as
motivation (Section~\ref{sec:results-motivation}).

\subsection{Motivation}
\label{sec:results-motivation}

Unmodified GMR output on floor-contact clips shows 12.9--15.9~cm maximum
floor penetration affecting 39--91\% of frames, against 1.0~cm and 0.3\% of
frames on GMR's own clean-locomotion clips (\texttt{walk1\_subject1}). This
gap holds at full-corpus scale, not only on individual clips.

\subsection{GMR's Own Mitigation is Insufficient}
\label{sec:results-heightfix}

\texttt{gmr\_heightfix} reduces penetration but leaves the robot floating
well above the floor at exactly the frames that should show tight ground
contact, since one rigid per-clip offset cannot satisfy a standing phase
and a lying phase in the same clip simultaneously
(Table~\ref{tab:heightfix}, \texttt{joint\_ok\%}).

\begin{table}[t]
\centering
\caption{\texttt{gmr\_heightfix}, class means.}
\label{tab:heightfix}
\begin{tabular}{lcc}
\toprule
Metric & Floor ($n=34$) & Locomotion ($n=43$) \\
\midrule
Floor penetration, max (cm)     & 2.76  & 3.06  \\
Floor penetration (\% frames)   & 0.38  & 3.42  \\
Self-collision (\% frames)      & 6.34  & 3.85  \\
Worst float (cm)                & 18.04 & 6.61  \\
Velocity spikes (mean/clip)     & 0.18  & 0.00  \\
Peak joint velocity (rad/s)     & 34.04 & 32.82 \\
Joint contact-correctness (\%)  & 0.19  & 46.26 \\
\bottomrule
\end{tabular}
\end{table}

\subsection{Multi-Stage Pipeline}
\label{sec:results-multistage}

\begin{table}[t]
\centering
\caption{Multi-stage pipeline, class means.}
\label{tab:multistage}
\begin{tabular}{lcc}
\toprule
Metric & Floor ($n=34$) & Locomotion ($n=43$) \\
\midrule
Floor penetration, max (cm)     & 0.00  & 0.00  \\
Floor penetration (\% frames)   & 0.00  & 0.00  \\
Self-collision (\% frames)      & 0.002 & 0.001 \\
Worst float (cm)                & 7.55  & 5.57  \\
Velocity spikes (mean/clip)     & 0.00  & 0.00  \\
Peak joint velocity (rad/s)     & 37.39 & 37.92 \\
Joint contact-correctness (\%)  & 98.85 & 98.79 \\
\bottomrule
\end{tabular}
\end{table}

Five of six metrics improve on \texttt{gmr\_heightfix} by more than an
order of magnitude; joint contact-correctness rises from 0.2\% to 98.8\% on
the floor class. Peak joint velocity is 9.8\% (floor) and 15.5\%
(locomotion) above \texttt{gmr\_heightfix} --- a real cost, attributable to
branch-flip artifacts in the per-frame correction stage
(Section~\ref{sec:discussion}) and only partially mitigated by a null-space
continuity term in that stage; the mitigation itself trades a smaller
regression in jerk and contact-slip on part of the floor class.

\subsection{Oracle Comparison}
\label{sec:results-oracle}

A per-clip constant Z-shift, swept by grid search to jointly minimize
penetration and floating with no notion of per-frame contact timing,
outperforms the multi-stage pipeline on either single-axis metric in
isolation (held-frame float within tolerance, or whole-body penetration,
each considered alone) on some clips. It cannot satisfy joint
contact-correctness, which requires both at the same frame: a rigid shift
cannot correct a multi-phase clip's standing stance and its lying phase
with one number. This is the basis for reporting joint contact-correctness
as the primary contact metric.

\subsection{Global Pipeline}
\label{sec:results-global}

Full-corpus evaluation of the global pipeline (Section~\ref{sec:global}) is
in progress; we report results on the 36 clips completed at time of writing
(18 floor-contact, 18 locomotion, class-balanced), pending completion of the
remaining 41.

\begin{table*}[t]
\centering
\caption{Global pipeline vs.\ \texttt{gmr\_heightfix} and the multi-stage
pipeline, class means, 36/77 clips.}
\label{tab:global}
\begin{tabular}{llccc}
\toprule
Metric & Class & \texttt{gmr\_heightfix} & Multi-stage & Global \\
\midrule
\multirow{2}{*}{Joint contact-correctness (\%)}
  & Floor      & 0.20  & 98.50 & 81.02 \\
  & Locomotion & 46.51 & 99.84 & 81.27 \\
\multirow{2}{*}{Floor penetration (cm)}
  & Floor      & 3.06  & 0.00  & 0.00  \\
  & Locomotion & 3.07  & 0.00  & 0.00  \\
\multirow{2}{*}{Self-collision (\% frames)}
  & Floor      & 6.14  & 0.01  & 1.13  \\
  & Locomotion & 2.87  & 0.00  & 0.91  \\
\multirow{2}{*}{Worst float (cm)}
  & Floor      & 18.35 & 8.16  & 13.54 \\
  & Locomotion & 7.87  & 5.03  & 11.53 \\
\multirow{2}{*}{Joint jerk}
  & Floor      & 5293  & 2601  & 803   \\
  & Locomotion & 7409  & 3101  & 1256  \\
\multirow{2}{*}{Peak joint velocity (rad/s)}
  & Floor      & 33.04 & 31.31 & 13.01 \\
  & Locomotion & 40.11 & 37.32 & 11.67 \\
\multirow{2}{*}{Velocity spikes (mean/clip)}
  & Floor      & 0.06  & 0.00  & 0.00  \\
  & Locomotion & 0.00  & 0.00  & 0.00  \\
\bottomrule
\end{tabular}
\end{table*}

The global pipeline improves on \texttt{gmr\_heightfix} on joint
contact-correctness, floor penetration, self-collision, and velocity spikes
on both classes, and reduces jerk and peak joint velocity relative to
\textit{both} baselines by a wide margin (jerk 3.2--6.6x lower, peak
velocity 2.4--3.4x lower than the multi-stage pipeline). It trails the
multi-stage pipeline on joint contact-correctness (81\% vs.\ 98--99\%) and
self-collision, consistent with a smaller total per-frame solve budget
(Section~\ref{sec:comparing}). Worst float is a mixed result: better than
\texttt{gmr\_heightfix} on the floor class (18.35 $\to$ 13.54~cm) but worse
on locomotion (7.87 $\to$ 11.53~cm), the one metric where the global
pipeline does not uniformly improve on the naive baseline. Which pipeline
is preferable depends on which axis matters more for a given downstream
use; Section~\ref{sec:disc-tradeoff} and the planned policy-training
evaluation (Section~\ref{sec:status-next-steps}) are intended to make that
trade-off decidable rather than qualitative.

\input{sections/discussion}
\input{sections/conclusions}

\bibliographystyle{IEEEtran}
\bibliography{../references}

\end{document}
